% Generated by scripts/assemble.py; corrections belong in source/corrections.json.
\examhead{164}{广西大学 2026 年数学分析试卷}{156}
% p157-b11; source PDF page 157
\par\Needspace{5\baselineskip}
\blockanchor{p157-b11}
\majorquestion{一}简答题.\par
% p157-b12; source PDF page 157
\blockanchor{p157-b12}
\question{1}已知点集 $S \subset \mathbb{R}$, 请叙述聚点的定义, 若 $S$ 为有限点集, 它是否存在聚点, 并给出理由.\par
% p157-b13; source PDF page 157
\blockanchor{p157-b13}
\question{2}已知 $I$ 为函数 $f(x)$ 定义域中的一个开区间, 请用 $\varepsilon - \delta$ 语言叙述 $f$ 在 $I$ 内不连续的定义, 并给出所有间断点的类型及定义.\par
% p157-b14; source PDF page 157
\blockanchor{p157-b14}
\question{3}请叙述二元函数 $f(x, y)$ 在点 $P_0(x_0, y_0)$ 处可微的定义及在该点的全微分，令 $z_0 = f(x_0, y_0)$，请写出 $f(x, y)$ 在点 $P_0(x_0, y_0)$ 处可微时，曲面 $z = f(x, y)$ 在点 $P(x_0, y_0, z_0)$ 处的切平面方程.\par
% p157-b15; source PDF page 157
\par\Needspace{5\baselineskip}
\blockanchor{p157-b15}
\majorquestion{二}计算题. 每题 15 分, 共 30 分.\par
% p157-b16; source PDF page 157
\blockanchor{p157-b16}
\question{1}设  $\left\{\begin{aligned}x(t)&=\int_{0}^{t}(e^{u}-1)\mathrm{d}u,\\ y(t)&=\ln(1+t^{2}).\end{aligned}\right.$  求  $\lim_{t\to0}\frac{y\sin y}{x^{2}}.$\par
% p157-b17; source PDF page 157
\blockanchor{p157-b17}
\question{2}求定积分  $I_{n}=\int_{0}^{\frac{\pi}{2}}\cos^{n}x\,\mathrm{d}x\;(n=1,2,\cdots)$ .\par
% p157-b18; source PDF page 157
\par\Needspace{5\baselineskip}
\blockanchor{p157-b18}
\majorquestion{三}证明题. 每题15分, 共90分.\par
% p157-b19; source PDF page 157
\blockanchor{p157-b19}
\question{1}设  $a_{1}, a_{2}, \cdots, a_{m}$  为 m 个正数，证明： $\sqrt[n]{a_{1}^{n} + a_{2}^{n} + \cdots + a_{m}^{n}}$  极限存在，并求极限值.\par
% p157-b20; source PDF page 157
\blockanchor{p157-b20}
\question{2}请叙述有限覆盖定理并用其证明: 若函数 $f(x)$ 在 $[a, b]$ 上连续, 则 $f(x)$ 在 $[a, b]$ 上一致连续.\par
% p157-b21; source PDF page 157
\blockanchor{p157-b21}
\question{3}请叙述罗尔定理和拉格朗日中值定理, 已知函数 $f(x)$ 在区间 $[a, b]$ 上满足拉格朗日中值定理的条件, $S(x)$ 表示以 $(a, f(a)), (b, f(b)), (x, f(x))$ 三点为顶点构成的三角形的面积, 请利用 $S(x)$ 和罗尔定理证明拉格朗日中值定理.\par
% p157-b22; source PDF page 157
\blockanchor{p157-b22}
\question{4}讨论函数 $\Gamma(p) = \int_{0}^{+\infty} x^{p-1} e^{-x} \, \mathrm{d}x$ 的收敛性并证明.\par
% p157-b23; source PDF page 157
\blockanchor{p157-b23}
\question{5}设函数 $f(x, y) \in C^2(\mathbb{R}^2)$，对任意的 $(x, y) \in \mathbb{R}^2$，都有 $f_{xx}(x, y) + f_{yy}(x, y) < 0$，请写出 $f(x, y)$ 的黑塞(Hesse)矩阵，并证明：$f(x, y)$ 在 $\mathbb{R}^2$ 上没有极小值点.\par
% p157-b24; source PDF page 157
\blockanchor{p157-b24}
\question{6}证明: 级数 $\frac{1!}{1} + \frac{2!}{2^2} + \frac{3!}{3^3} + \cdots + \frac{n!}{n^n} + \cdots$ 收敛.\par
